% Part: history
% Chapter: biographies
% Section: kurt-goedel

\documentclass[../../../include/open-logic-section]{subfiles}

\begin{document}

\olfileid{his}{bio}{god}

\olsection{کورت گودل}

\olphoto{goedel-kurt}{Kurt G{\"o}del}

کورت گودل (\textsc{ger}-dle) در 28 آوریل 1906 در برونِ امپراتوری
اتریش--مجارستان، شهر برنوی کنونی در جمهوری چک، به دنیا آمد. خانواده‌اش
به سبب طبیعت کنجکاو و هوش او، کورتلهٔ جوان را اغلب «Der kleine Herr
Warum» یعنی «آقای کوچک چرا» می‌نامیدند. او از دبستان به بعد در درس‌ها
درخشان بود و فقط در ریاضیات نمره‌ای پایین‌تر از بالاترین نمره گرفت.
گودل به سبب سلامتی نامناسب اغلب از مدرسه غایب بود و از تربیت بدنی معاف
شد. در کودکی ابتلای او به تب روماتیسمی تشخیص داده شد. او در سراسر
زندگی، برخلاف ارزیابی پزشکان، باور داشت این بیماری برای همیشه به قلبش
آسیب زده است.

گودل در سال 1924 تحصیل در دانشگاه وین را آغاز کرد و دورهٔ دکتری را در
~1929 به پایان رساند. ابتدا قصد داشت فیزیک بخواند، اما علاقه‌اش به‌زودی
به ریاضیات و به‌ویژه منطق معطوف شد؛ نفوذ رودلف کارناپِ فیلسوف یکی از
علت‌ها بود. رساله‌اش که زیر نظر هانس هان نوشته شد، قضیهٔ تمامیت منطق
محمول‌های مرتبهٔ اول با همانی را اثبات کرد \citep{Godel1929}. تنها یک
سال بعد، مشهورترین نتایج خود، یعنی قضیه‌های اول و دوم ناتمامیت، را به
دست آورد که در \citealt{Godel1931} منتشر شدند. گودل در دوران اقامت در
وین فعالیت گسترده‌ای در حلقهٔ وین داشت؛ گروهی از فیلسوفان با گرایش
علمی که کارناپ نیز عضو آن بود و نتایج گودل تأثیر ویژه‌ای بر کارش گذاشت.

گودل در سال 1938 با ادل نیمبورسکی ازدواج کرد. والدینش خشنود نبودند:
ادل نه‌تنها شش سال از او بزرگ‌تر و پیش‌تر طلاق گرفته بود، بلکه به‌عنوان
رقاص در باشگاهی شبانه کار می‌کرد. بااین‌حال فشارهای اجتماعی بر گودل
اثر نگذاشت و آن دو تا زمان مرگ او زندگی مشترک شادی داشتند.

پس از الحاق اتریش به آلمان نازی در 1938، گودل و ادل به ایالات متحده
مهاجرت کردند و او در مؤسسهٔ مطالعات پیشرفته در پرینستون نیوجرسی
جایگاهی گرفت. با وجود درون‌گرایی و منش نامتعارفش، دوران گودل در
پرینستون با همکاری و ثمربخشی همراه بود. او جستارهایی در نظریهٔ
مجموعه‌ها، فلسفه و فیزیک منتشر کرد. به‌ویژه دوستی بسیار نزدیکی با
آلبرت اینشتین، همکارش در مؤسسه، برقرار کرد.

در سال‌های پایانی زندگی، سلامت روان گودل رو به وخامت گذاشت. بستری‌شدن
همسرش در سال 1977 به این معنا بود که دیگر نمی‌توانست برای او غذا
بپزد. گودل که در سراسر زندگی با مشکلات سلامت روان روبه‌رو بود، گرفتار
بدگمانی شدید شد. از ترس مرگبار مسموم‌شدن از خوردن غذا سر باز زد و در
14 ژانویهٔ 1978 در پرینستون بر اثر گرسنگی درگذشت.

\begin{reading}
برای زندگی‌نامه‌ای کامل از گودل، \citet{Dawson1997} را ببینید. برای
مطالب زندگی‌نامه‌ای بیشتر و نیز جستارهایی دربارهٔ سهم گودل در منطق و
فلسفه، \citet{Wang1990}، \citet{Baaz2011}، \citet{Takeuti2003} و
\citet{Sigmund2007} را ببینید.

رسالهٔ دکتری گودل به زبان اصلی آلمانی در دسترس است
\citep{Godel1929}. متن اصلی قضیه‌های ناتمامیت در \citep{Godel1931}
آمده است. همهٔ نوشته‌های منتشرشده و منتشرنشدهٔ گودل، همراه با گزیده‌ای
از نامه‌های او، به انگلیسی در \emph{Collected Papers} او در دسترس‌اند
\citet{Godel1986,Godel1990}.

برای بررسی مفصل قضیه‌های ناتمامیت گودل، \citet{Smith2013} را ببینید.
برای بحثی غیررسمی و فلسفی دربارهٔ قضیه‌های گودل نیز به پادکست مارک
لینسن‌مایر مراجعه کنید \citep{Linsenmayer2014}.
\end{reading}

\end{document}
